\section{CƠ SỞ LÝ THUYẾT}

\subsection{Metro Ethernet MAN}

Metro Ethernet là mô hình cung cấp dịch vụ Ethernet trong phạm vi đô thị hoặc khu vực rộng hơn LAN truyền thống. Đối với doanh nghiệp nhiều chi nhánh, Metro Ethernet cho phép các site ở xa nhau kết nối như một hạ tầng thống nhất qua mạng của nhà cung cấp. Ưu điểm chính là tính quen thuộc của Ethernet, khả năng mở rộng băng thông và khả năng tích hợp với các dịch vụ VPN/MPLS của provider \cite{mef}.

\subsection{Vai trò CE, PE và P router}

Trong mô hình provider:
\begin{itemize}
    \item CE (Customer Edge) nằm phía khách hàng, kết nối LAN chi nhánh vào mạng nhà cung cấp.
    \item PE (Provider Edge) nằm ở biên provider, nhận traffic từ CE và đưa traffic vào miền MPLS.
    \item P (Provider) nằm trong lõi provider, không cần biết chi tiết LAN khách hàng mà chủ yếu chuyển tiếp dựa trên label.
\end{itemize}

Phân tách CE/PE/P giúp báo cáo mô phỏng đúng vai trò của từng miền: khách hàng chịu trách nhiệm LAN nội bộ, provider chịu trách nhiệm backbone và chuyển tiếp liên chi nhánh.

\subsection{Nguyên lý MPLS label switching}

MPLS gắn một shim header 4 byte giữa Ethernet header và IP packet. Router MPLS xử lý gói dựa trên label:
\begin{itemize}
    \item Push label: ingress PE thêm label vào gói IP khi đi vào miền MPLS.
    \item Swap label: P router thay label đầu vào bằng label đầu ra tương ứng.
    \item Pop label: egress PE bỏ label trước khi chuyển gói về CE và LAN đích.
\end{itemize}

Trong project, quá trình này được chứng minh bằng hai nhóm log. Bảng route MPLS cho thấy các dòng \texttt{encap mpls}, \texttt{proto ldp} và \texttt{as to} trong backbone. Tcpdump trên core link bắt được ethertype \texttt{0x8847} và label động do LDP cấp phát, chứng minh traffic đi qua miền MPLS bằng label thật.

\subsection{OSPF, LDP, VPLS và cách project xử lý yêu cầu rubric}

Rubric có nhắc đến OSPF, LDP và VPLS ở mức điểm tối đa. Trong mạng MPLS nhà cung cấp thực tế, OSPF/IS-IS thường dùng để hội tụ định tuyến lõi, LDP dùng để phân phối label động, còn VPLS cung cấp dịch vụ L2VPN đa điểm.

Project này triển khai MPLS native bằng Linux kernel và tích hợp FRRouting trong các namespace PE/P. OSPF dùng để học router-id và mạng lõi, LDP phân phối label động trong backbone; bằng chứng được lưu ở \texttt{results/frr\_control\_plane.txt} với trạng thái OSPF \texttt{Full/-}, LDP \texttt{OPERATIONAL}, bảng binding và \texttt{show mpls table}. Luồng CE/branch được đưa vào service VPLS/L2VPN; data-plane L2VPN trong Mininet dùng full-mesh Linux GRETAP pseudowire có split-horizon/port-isolation trên underlay MPLS để có forward thật và kiểm chứng được.

\subsection{Các kiến trúc LAN trong ba chi nhánh}

Branch 1 dùng Flat Network, trong đó host nằm cùng một subnet và đi qua ít thiết bị trung gian. Kiến trúc này đơn giản, dễ cấu hình nhưng khó mở rộng khi quy mô lớn.

Branch 2 dùng Three-tier Network gồm Access, Distribution và Core. Mô hình này phù hợp mạng doanh nghiệp truyền thống vì tách rõ chức năng từng lớp, thuận lợi cho mở rộng, quản trị và áp chính sách.

Branch 3 dùng Leaf-Spine Network, thường gặp trong data center. Host gắn vào leaf, leaf nối lên spine để tạo đường đi đồng đều hơn. Trong lab, topology được giữ loop-free ở L2 để Mininet/OVS khởi động ổn định.
